\documentclass[11pt]{article}
\usepackage[margin=25mm]{geometry}
\usepackage[T1]{fontenc}
\usepackage{booktabs,url,hyperref}
\hypersetup{colorlinks=true,urlcolor=blue}
\title{Research Artifact Preservation and Audio Publication Supplement}
\author{AI--Human Music Detection Thesis Project}
\date{13 September 2026: versioned delivery checkpoint}
\begin{document}
\maketitle

\section{Scope and scientific status}
This supplement records preservation work after the English thesis v7. It
does not change feature definitions, model fits, thresholds, experimental
memberships, or reported classification results. The BC extension, expanded
YuE Native30 evaluations, locked scoring, legacy spectral comparison and
Native60 transfer experiments are complete. Whole-project audio delivery is
\emph{not} complete. Publication status must not be substituted for scientific
validation, and file counts must not be interpreted as independent songs.

Code and versioned catalogues are maintained in
\href{https://github.com/ezmonyi/ai-music-detection}{the project GitHub repository}.
Audio is published in
\href{https://huggingface.co/datasets/EZMONYI/music-ai-human-test-audio}{the public HF dataset}.
The local report root, including older results, is
\nolinkurl{/Users/yi/Documents/report/music_ai_detection_20260912}.

\section{Original-source and actual-input coverage are different}
At the fixed HF revision
\nolinkurl{caed9220f68b258e9c2dcb89792ccb1fe4d1dc61}, the public snapshot
contained 16,795 direct audio/MIDI paths and 130,247,953,279 logical bytes.
Archive members were not enumerated. Of 10,141 historical recorded raw hashes,
7,781 matched public audio objects and 2,360 did not. These counts exclude
later-only identities and do not prove processed-input coverage.

A separate audit read the actual \texttt{audio\_path} files named in the
historical 10-second and 30-second metadata. All 14,246 distinct filesystem
paths were readable. Their 14,638 experimental memberships resolve to 14,244
unique byte objects. Only 966 objects matched the pre-supplement public
snapshot; 13,278 objects, totaling 19,083,227,708 bytes, did not. The larger
13,480 unmatched count from the initial audit refers to memberships, not
unique files. Runtime crop settings are retained; the audit does not claim
that every runtime crop was rendered or every separated stem was enumerated.

\section{Independently accepted processed-input supplements}
\begin{center}
\begin{tabular}{lrr}
\toprule
Supplement & File objects & Bytes \\
\midrule
DiffRhythm pilot & 100 & 410,987,527 \\
Ishizaka/MusicNet subset & 78 & 57,343,890 \\
Historical MAESTRO & 600 & 1,259,759,294 \\
MagnaTagATune & 500 & 218,959,824 \\
\midrule
Total & 1,278 & 1,947,050,535 \\
\bottomrule
\end{tabular}
\end{center}
Each acceptance checks an exact public manifest and notice at one fixed
revision and every remote file hash/size. It is not a full waveform
redownload. DiffRhythm retains 50 outputs and ten conditioning groups;
Ishizaka retains 39 source-reviewed recordings; MAESTRO retains 300
performances. These are processed versions of existing inputs.

The acceptance records reside under \texttt{current\_results/} in folders
\nolinkurl{hf_diffrhythm_test_views_v1},
\nolinkurl{hf_ishizaka_test_views_v1},
\nolinkurl{hf_maestro_historical_views_v1}, and
\nolinkurl{hf_mtt_test_views_v1}. Each contains
\nolinkurl{INDEPENDENT_ACCEPTANCE.json} and \nolinkurl{manifest.json}.
The GitHub dataset catalogues preserve corresponding receipts and batch records.

\section{Remaining processed-file backlog}
After these acceptances, 12,000 objects (17,136,177,173 bytes) remain without
final acceptance in the frozen audit scope. A dispatched publisher is not
proof of a live transfer or completed release.
\begin{center}
\begin{tabular}{lr}
\toprule
Checkpoint category & Objects \\
\midrule
AIME: dispatched, not accepted & 5,500 \\
ACE-Step/HeartMuLa: dispatched, not accepted & 800 \\
FMA licensed candidates: dispatched, not accepted & 316 \\
Suno: prepared, not dispatched & 900 \\
FMA: no-derivatives restriction & 390 \\
FMA: unresolved or mixed terms & 194 \\
Other source rights/provenance unresolved & 3,900 \\
\bottomrule
\end{tabular}
\end{center}
The FMA screening retains exact license versions and jurisdiction where
recorded; it does not relicense all files under a single modern CC license.
Metadata declarations are not a warranty of every underlying recording right.
ND processing restrictions are handled conservatively. Other source issues
include individual recording rights, creator attribution and explicit
redistribution-consent requests. Preserved bytes are not equivalent to
permission to publish them.

\section{Upload recovery and reproducibility safeguards}
Two distinct HF limits were observed: short-window API limits and a response
stating 128 repository commits per hour. The latter terminated the AIME
publisher after 1,900 processed files. Earlier short-cooldown recovery attempts
for open-model and FMA files also returned 429 responses. Verified process
identities were used when replacing those outdated recovery workers; existing
audio, batch receipts and logs were preserved. Replacement workers wait
3,900 seconds before resuming and honor longer server retry delays. This
checkpoint does not claim that the cooldown has ended or uploads have resumed.

Authentication is passed through standard input, not committed to the
repository. Recovery uses per-output process locks. Completed batches are
rechecked against remote hashes before being skipped; new files are checked
against their expected local hashes before upload. These checks are separate
from the final anonymous acceptance.

\section{Code and analysis references}
Current scripts are under the GitHub directory
\nolinkurl{current/yue2_500_extension_20260911/}. Relevant implementations are:
\begin{itemize}
\item \nolinkurl{audit_historical_test_views_v1.py} and\newline
\nolinkurl{plan_historical_test_view_objects_v1.py}: actual-input hashing and deduplication.
\item \nolinkurl{reconcile_processed_delivery_v1.py} and
\nolinkurl{classify_processed_backlog_v1.py}: receipt-bound coverage and backlog partitioning.
\item \nolinkurl{plan_fma_test_view_rights_v1.py}: conservative per-recording license screening.
\item \nolinkurl{resume_processed_upload_v1.py}: locked, rate-aware recovery.
\item Source-specific \texttt{publish\_*\_test\_views\_v1.py} and
\texttt{verify\_*\_test\_views\_v1.py}: publication and acceptance.
\end{itemize}
Analysis products are preserved under local \texttt{datasets/}:
\nolinkurl{historical_public_coverage_20260913_v2},
\nolinkurl{historical_test_view_objects_v1},
\nolinkurl{processed_delivery_reconciliation_v1},
\nolinkurl{processed_backlog_v1}, and \nolinkurl{fma_test_view_rights_v2}.
Their COMMIT records bind output hashes and exact audit scopes.

Six upload-backoff tests and four license-screening tests passed locally.
Byte comparisons found all 162 current YuE-extension Python scripts and all
392 current audio-phenomena Python scripts matched GitHub staging. A separate
1,223-file publication manifest records source hashes and zero matches for
the checker's literal credential patterns. Neither check constitutes a full
security audit, complete environment capture, or all-project source coverage.

\section{Outstanding work}
Finish and independently accept remaining permitted uploads, dispatch the
prepared Suno publication when quota allows, reconcile unenumerated stems and
earlier/external scopes, resolve remaining source permissions, and update
the final delivery index and thesis references without overwriting this
checkpoint. No experimental result is changed by this supplement.
\end{document}
